\section{Resultados}

\subsection{Desempenho no treino}

A Tabela~\ref{tab:treino} resume as métricas no conjunto de treino. O filtro
casado e a combinação Wiener\,+\,casado obtêm AUC \(\valAucCasadoTrain\),
superior ao limiar baseado no máximo de \(|y|\) (sinal bruto ou Wiener). O RMSE
da forma cai de cerca de \(0{,}50\) (bruto/Wiener) para cerca de \(0{,}13\)
quando se reconstrói \(\hat{A}\,s[n-\hat{n}_0]\); o BLUE colorido reduz ainda
um pouco esse erro e melhora o RMSE de amplitude.

O limiar a FA \(\approx 1\%\) é conservador: a precisão é alta, mas a
revocação é baixa, o que reduz o F1 (por exemplo, \(\valFoneCasadoTrain\) para o
casado). A matriz de confusão do melhor detector no treino
(Figura~\ref{fig:cm-train}) ilustra esse compromisso: poucos falsos positivos
(\(\valFpTrain\)) e muitos falsos negativos (\(\valFnTrain\)).

\begin{table}[H]
\centering
\caption{Métricas no conjunto de treino.}
\label{tab:treino}
\resizebox{\textwidth}{!}{\begin{tabular}{lcccccccc}
\toprule
Método & AUC & Prec. & Rev. & F1 & RMSE forma & RMSE $A$ & med$|\Delta n_0|$ & FA emp. \\
\midrule
Sinal bruto & 0{,}675 & 0{,}952 & 0{,}130 & 0{,}229 & 0{,}496 & 1{,}117 & 11{,}00 & 0{,}010 \\
Wiener & 0{,}702 & 0{,}980 & 0{,}160 & 0{,}275 & 0{,}139 & 2{,}381 & 7{,}00 & 0{,}005 \\
Casado + BLUE (branco) & 0{,}731 & 0{,}985 & 0{,}208 & 0{,}344 & 0{,}129 & 1{,}020 & 4{,}00 & 0{,}005 \\
Wiener + casado & 0{,}722 & 0{,}979 & 0{,}153 & 0{,}265 & 0{,}134 & 1{,}008 & 5{,}00 & 0{,}005 \\
Casado + BLUE (colorido) & 0{,}731 & 0{,}985 & 0{,}208 & 0{,}344 & 0{,}121 & 0{,}856 & 4{,}00 & 0{,}005 \\
\bottomrule
\end{tabular}
}
\end{table}

\begin{figure}[H]
\centering
\includegraphics[width=0.75\textwidth]{08_roc_treino.png}
\caption{Curvas ROC no conjunto de treino.}
\label{fig:roc-train}
\end{figure}

\begin{figure}[H]
\centering
\includegraphics[width=0.55\textwidth]{10_confusao_treino.png}
\caption{Matriz de confusão no treino (detector escolhido).}
\label{fig:cm-train}
\end{figure}

\begin{figure}[H]
\centering
\includegraphics[width=0.9\textwidth]{12_reconstrucao_exemplo.png}
\caption{Exemplo de reconstrução do pulso a partir do BLUE e comparação com Wiener.}
\label{fig:recon}
\end{figure}

\subsection{Desempenho no teste}

Com limiares e filtros congelados, as métricas no teste (Tabela~\ref{tab:teste})
acompanham o comportamento do treino. A AUC do casado sobe ligeiramente para
\(\valAucCasadoTest\), enquanto o F1 cai para \(\valFoneCasadoTest\) --- de novo
pelo limiar conservador. O RMSE da forma permanece próximo
(\(\valRmseWaveCasadoTrain\) no treino e \(\valRmseWaveCasadoTest\) no teste),
e o BLUE colorido mantém o melhor RMSE de amplitude entre os métodos
comparados.

\begin{table}[H]
\centering
\caption{Métricas no conjunto de teste.}
\label{tab:teste}
\resizebox{\textwidth}{!}{\begin{tabular}{lcccccc}
\toprule
Método & AUC & F1 & RMSE forma & RMSE $A$ & med$|\Delta n_0|$ & FA emp. \\
\midrule
Sinal bruto & 0.661 & 0.154 & 0.494 & 0.982 & 3.00 & 0.033 \\
Wiener & 0.661 & 0.154 & 0.493 & 0.982 & 3.00 & 0.033 \\
Filtro casado & 0.756 & 0.233 & 0.122 & 0.982 & 3.00 & 0.008 \\
Wiener + casado & 0.756 & 0.233 & 0.122 & 0.982 & 3.00 & 0.008 \\
BLUE (branco) & 0.756 & 0.233 & 0.122 & 0.982 & 3.00 & 0.008 \\
BLUE (colorido) & 0.756 & 0.233 & 0.112 & 0.843 & 3.00 & 0.008 \\
\bottomrule
\end{tabular}
}
\end{table}

\begin{figure}[H]
\centering
\includegraphics[width=0.75\textwidth]{09_roc_teste.png}
\caption{Curvas ROC no conjunto de teste.}
\label{fig:roc-test}
\end{figure}

\begin{figure}[H]
\centering
\includegraphics[width=0.55\textwidth]{11_confusao_teste.png}
\caption{Matriz de confusão no teste (mesmo limiar do treino).}
\label{fig:cm-test}
\end{figure}

\subsection{Comparação dos métodos}

Com base apenas no treino, e confirmado no teste:

\begin{itemize}
  \item \textbf{Reconstrução da forma:} melhor com \valBestWave, isto é,
  estimando \(\hat{A}\) e reposicionando o template; o sinal bruto e a saída de
  Wiener isolada ficam bem atrás em RMSE.
  \item \textbf{Detecção:} o filtro casado (e Wiener\,+\,casado, com desempenho
  praticamente idêntico) supera o uso do máximo de \(|y|\). A pré-filtragem de
  Wiener não melhorou de forma relevante a AUC neste conjunto.
  \item \textbf{Amplitude:} \valBestAmp reduz o RMSE em relação ao BLUE branco
  e ao baseline do máximo da observação, o que é coerente com o caráter
  colorido do ruído.
\end{itemize}
